\section{The actual tangent-displacement-zero quotient}

Write $B=\mathcal O_{\mathbb C^d,0}$, $R=\mathcal O_{\mathbb C^{d+t},0}$
and $S=\mathcal O_{\mathbb C^{d+(d+t)},0}$. The coordinates of $S$ are
$(c;v,x)$, where $c,v\in\mathbb C^d$ and $x\in\mathbb C^t$.
The existing joint map is the actual complex linear map
$L(c;v,x)=(c+v,x)$, and both $B$-actions are the original pullbacks of
the $c$ coordinates. All ideals below are whole ideals of holomorphic
germs. This chapter supplies a coordinate quotient; it asserts neither
normal flatness nor the remaining general Bennett inequality.

\begin{definition}[The literal zero-displacement trace]
\label{hir:tangent-displacement-zero-trace}
\lean{Hironaka.HironakaActualTangentDisplacementZeroQuotient.actualTangentDisplacementZeroInclusion,Hironaka.HironakaActualTangentDisplacementZeroQuotient.actualTangentDisplacementZeroTrace,Hironaka.HironakaActualTangentDisplacementZeroQuotient.tangentDisplacementIdeal}
\leanok
Let $E:\mathbb C^{d+t}\longrightarrow\mathbb C^{d+(d+t)}$ be
$E(c,x)=(c;0,x)$. Let $q=E^*:S\longrightarrow R$ and let
$V=(v_0,\ldots,v_{d-1})\subset S$.
\end{definition}

\begin{lemma}[Actual joint and scalar retractions]
\label{hir:tangent-displacement-zero-retractions}
\uses{hir:tangent-displacement-zero-trace}
\lean{Hironaka.HironakaActualTangentDisplacementZeroQuotient.actualTangentDisplacementZeroTrace_jointPullback,Hironaka.HironakaActualTangentDisplacementZeroQuotient.actualTangentDisplacementZeroTrace_tangentPullback}
\leanok
For every $f\in R$ and $b\in B$, one has $q(L^*f)=f$ and
$q(p_S^*b)=p_R^*b$, where $p_S(c;v,x)=c$ and $p_R(c,x)=c$.
\end{lemma}
\begin{proof}
The actual coordinate formulas give $L(E(c,x))=(c,x)$ and
$p_S(E(c,x))=c=p_R(c,x)$ at every point. For a holomorphic
representative $F$ of $f$, the representative of $q(L^*f)$ is
$F\circ L\circ E=F$. Equality on the entire model implies equality
of germs. Applying the same argument to a representative of $b$
proves the second equality. In particular, $q$ respects the specified
original $B$-actions.
\end{proof}

\begin{theorem}[The entire trace kernel is the original middle coordinate ideal]
\label{hir:tangent-displacement-zero-kernel}
\uses{hir:tangent-displacement-zero-retractions,hir:relative-normal-entire-plane-kernel}
\lean{Hironaka.HironakaActualTangentDisplacementZeroQuotient.actualTangentDisplacementZeroTrace_surjective,Hironaka.HironakaActualTangentDisplacementZeroQuotient.tangentDisplacementIdeal_eq_ker}
\leanok
The genuine map $q:S\longrightarrow R$ is surjective and $\ker q=V$.
\end{theorem}
\begin{proof}
The preceding retraction supplies the preimage $L^*f$ of every $f\in R$.
Permute the three coordinate blocks by $(c;v,x)\mapsto(c,x;v)$.
This is a bijective complex linear coordinate map with its literal inverse,
so its pullback $\rho$ is an isomorphism of whole holomorphic germ algebras.
Under this permutation, $q$ is the existing trace which kills the last
$d$ coordinates while retaining the first $d+t$ coordinates. That trace
has kernel the entire ideal of its last $d$ coordinate germs by the
coordinate-plane Hadamard theorem. Moreover $\rho(V)$ is exactly that
ideal: its actual generators are the same last $d$ coordinates. Pulling
back the kernel equality through the bijection $\rho$ proves
$\ker q=V$. This argument retains all germs in the kernel, without
restricting to a chosen family of generators of an unrelated ideal.
\end{proof}

\begin{theorem}[The same whole original ideal is recovered]
\label{hir:tangent-displacement-zero-original-ideal}
\uses{hir:tangent-displacement-zero-retractions}
\lean{Hironaka.HironakaActualTangentDisplacementZeroQuotient.map_jointIdeal_eq_original,Hironaka.HironakaActualTangentDisplacementZeroQuotient.map_actualCentreJointIdeal_eq_original}
\leanok
For every ideal $I\subset R$, put $J=L^*I$, meaning ideal extension
along the actual whole joint pullback. Then $q_*J=I$.
For the existing joint ideal generated by all $L^*H_i$, its image is
literally the ideal generated by all original germs $H_i$.
\end{theorem}
\begin{proof}
Extensions of ideals along two homomorphisms compose. Their composite
here is $q\circ L^*=\mathrm{id}_R$ by the retraction already proved.
Hence $q_*(L^*I)=I$ for every whole ideal $I$.
For the second assertion, ideal extension sends the span of a set to
the span of its image. Each generator $L^*H_i$ maps to the same $H_i$,
so the two displayed spans coincide, with no additional hypothesis on
the indexing family.
\end{proof}

\begin{theorem}[The genuine displacement ideal maps to the original normal ideal]
\label{hir:tangent-displacement-zero-normal-ideal}
\uses{hir:tangent-displacement-zero-trace}
\lean{Hironaka.HironakaActualTangentDisplacementZeroQuotient.map_jointNormalIdeal_eq_normal,Hironaka.HironakaActualTangentDisplacementZeroQuotient.map_jointNormalIdeal_pow_eq_normal}
\leanok
Let $K=(v_0,\ldots,v_{d-1},x_0,\ldots,x_{t-1})\subset S$ be
the existing original joint relative normal ideal and let
$N=(x_0,\ldots,x_{t-1})\subset R$ be the original normal ideal.
Then $q_*K=N$, and $q_*(K^a)=N^a$ for every $a\geq0$.
\end{theorem}
\begin{proof}
For every $i<d$, the representative of $q(v_i)$ is the identically
zero coordinate function, so $q(v_i)=0$. For every $j<t$, the
representative of $q(x_j)$ is the same actual coordinate function
$x_j$ on the original model, so $q(x_j)=x_j$. Thus the image ideal
is contained in $N$. Conversely, each generator $x_j$ of $N$ is the
image of the corresponding genuine generator of $K$, proving the
reverse inclusion. Ideal extension preserves products, and induction
on $a$ now gives $q_*(K^a)=N^a$, including $a=0$.
\end{proof}

\begin{corollary}[The same finite truncation ideals have their original images]
\label{hir:tangent-displacement-zero-truncated-ideal}
\uses{hir:tangent-displacement-zero-original-ideal,hir:tangent-displacement-zero-normal-ideal}
\lean{Hironaka.HironakaActualTangentDisplacementZeroQuotient.map_jointTruncation_eq_original}
\leanok
For every original ideal $I\subset R$ and every $a\geq0$, one has
$q_*(L^*I+K^a)=I+N^a$.
\end{corollary}
\begin{proof}
Ideal extension preserves sums. Apply the proved whole-original-ideal
image equality to the first summand and the proved normal-power image
equality to the second summand. Their sum is exactly $I+N^a$.
\end{proof}

\begin{theorem}[The original tangent scalar quotient algebra]
\label{hir:tangent-displacement-zero-base-algebra}
\uses{hir:tangent-displacement-zero-kernel,hir:tangent-displacement-zero-retractions}
\lean{Hironaka.HironakaActualTangentDisplacementZeroQuotient.actualTangentDisplacementZeroBaseAlgHom,Hironaka.HironakaActualTangentDisplacementZeroQuotient.tangentDisplacementQuotientAlgEquiv,Hironaka.HironakaActualTangentDisplacementZeroQuotient.tangentDisplacementQuotientAlgEquiv_mk}
\leanok
The actual trace is a $B$-algebra homomorphism for the original tangent
pullback actions and induces an isomorphism $S/V\simeq_B R$.
\end{theorem}
\begin{proof}
The scalar retraction identifies the image of every $b\in B$ with its
original scalar in $R$. Together with the ring homomorphism laws of
pullback this is exactly the $B$-algebra homomorphism property.
The map is surjective and has kernel $V$, as proved above. The first
isomorphism theorem for this same $B$-algebra homomorphism gives the
displayed isomorphism, with every class represented by $s\in S$
mapping to its literal trace $q(s)$. No coordinate quotient or scalar
compatibility is assumed as an input.
\end{proof}
